\begin{helper}[Third exit when the supplied exits lie in distinct components]
Let $G=(V,E)$ be a finite loopless multigraph whose elements of $E$ are
exact edge identities, let $T$ be a tree-sequence, and let $\varphi$ be a
proper partial $k$-edge-coloring.  Suppose that $V(T)$ is elementary.
Let $a,b\in[k]$ be distinct, with $a$ missing on $T$, the singleton
$\{a\}$ closed on $T$, and $b$ not missing on $T$.  Let $P_1,P_2$ be
support-disjoint $T$-exit paths with outside ends $a_1\ne a_2$ and tree
ends $b_1\ne b_2$.  If the full $(a,b)$-components through $a_1$ and
$a_2$ are distinct, then there are a third outside end $a_3$, a third
tree end $b_3$, a $T$-exit path $P_3$ from $a_3$ to $b_3$, and an exact
edge $f_3$ such that $P_3$ is support-disjoint from both supplied paths,
its outside and tree ends differ respectively from both preceding ends,
and $f_3$ lies on $P_3$, is incident with $b_3$, and is a $b$-colored
boundary edge of $V(T)$.
\end{helper}
\begin{proof}
Put $X=V(T)$ by \noderef{TreeSequenceVertexSet}, and let $C_i$ be the
full \noderef{TwoColorComponent} through the vertex-node $a_i$.  We first
record the full-component geometry forced by either supplied exit path.
The clauses of \noderef{TreeExitPath} say that $a_i\notin X$, that at
least one of $a,b$ is missing at $a_i$, and that $P_i$ has length at
least two.  Thus $P_i$ supplies a neighbor of $a_i$ in
\noderef{TwoColorIncidenceGraph}.  If $a$ is missing at $a_i$, every such
neighbor is the identity of a $b$-colored edge incident with $a_i$; two
distinct neighbors would be distinct adjacent edges of the same color,
contrary to \noderef{AdjacentEdges} and
\noderef{IsProperEdgeColoring}.  The same argument with the colors
reversed applies when $b$ is missing.  Hence $a_i$ has exactly one
neighbor in its full component.

The path $P_i$ supplies an edge-node in $C_i$, while this degree-one
vertex-node supplies the leaf in \noderef{TwoColorPathComponent}; hence
$C_i$ is a \noderef{TwoColorNontrivialPathComponent}.  Apply
\noderef{TwoColorComponentSpanningPath}.  Its spanning walk is simple and
contains every node of $C_i$.  The degree-one node $a_i$ must be an end
of that path, since an internal occurrence has two distinct path
neighbors.  By
\noderef{TwoColorIncidenceDegreeAtMostTwo}, no internal node of the
spanning path has an additional neighbor.  Nor can an end have an
additional neighbor: a neighbor at a nonadjacent internal position would
give degree at least three there, while an adjacency joining the two ends
would give $a_i$ a second neighbor.  Consequently the full component, including all
its adjacencies, is exactly a finite nontrivial simple path, with $a_i$ as
one end.

The other end is a vertex-node, say $c_i$.  Indeed, an edge-node of the
two-color incidence graph represents a colored exact edge.  Its two
endpoint entries are distinct by \noderef{LooplessMultigraph}, and the
definition of \noderef{Incident} makes their vertex-nodes two distinct
neighbors, so an edge-node cannot be an end.  Moreover $c_i$ misses $a$
or $b$: if both colors were present there, \noderef{PresentColors} would
supply distinct incident edges of the two colors and hence two distinct
neighbors.  By the complementary definition \noderef{MissingColors},
this proves the claimed missing-color alternative.

Since $C_1\ne C_2$, \noderef{TwoColorComponentsDisjoint} makes the two
full components disjoint.  Thus their four ends $a_1,c_1,a_2,c_2$ are
pairwise distinct.  No vertex of $X$ misses $b$, because otherwise the
union definition \noderef{TreeMissingColors} would put $b$ among the
tree-missing colors.  Therefore, if $c_i\in X$, its endpoint
missing-color alternative forces it to miss $a$.  At most one of
$c_1,c_2$ can do so: two distinct vertices of $X$ missing $a$ contradict
\noderef{ColoredVertexSetElementary}.  Hence at least one of $c_1,c_2$
lies outside $X$.  Interchanging the labels $1,2$ in the following
symmetric argument if necessary, assume $c_1\notin X$.

Orient the full simple path $C_1$ from $c_1$ toward $a_1$, and let $b_3$
be the first vertex on it which lies in $X$; such a vertex exists because
$b_1\in X$ lies on the supplied path and hence in $C_1$.  Let $P_3$ be
the initial portion ending at $b_3$, and put $a_3=c_1$.  This portion is
simple, begins outside $X$, ends inside $X$, and has no other vertex of
$X$.  Its first vertex misses $a$ or $b$ by the endpoint analysis above.
Its endpoints are distinct vertex-nodes, and the incidence graph is
bipartite between vertex-nodes and edge-nodes, so its length is at least
two.  These are exactly all clauses of \noderef{TreeExitPath}.

The component disjointness already makes $P_3$ support-disjoint from
$P_2$.  It remains to compare $P_3$ with $P_1$ inside $C_1$.  Orient the
full component path from $a_1$ to $c_1$.  The supplied path $P_1$ is its
initial segment ending at $b_1$: starting at the degree-one node $a_1$,
its first edge is forced; inductively, at every internal node the next
edge is the unique component-path edge different from the preceding one,
because there is no additional component adjacency.  Simplicity prevents
backtracking.  Its endpoint is its first vertex in $X$ by
\noderef{TreeExitPath}.  Similarly, $P_3$ is the terminal segment from
$c_1$ through the first vertex of $X$ encountered in the reverse order.

If these two segments met, their path-index intervals would force the
reverse first contact $b_3$ to occur no later than the forward first
contact $b_1$.  Then $b_3$ lies on $P_1$; the tree-intersection clause of
\noderef{TreeExitPath} forces $b_3=b_1$.  Let $g_1,g_3$ be the exact
edge-nodes immediately preceding this common vertex on $P_1,P_3$.
They are distinct, since they occur on opposite sides of the common
vertex in a simple path whose two initial ends differ.  The vertex before
each $g_i$ lies outside $X$, while the common endpoint lies in $X$.
Using \noderef{TwoColorIncidenceGraph}, \noderef{Incident}, and the
loopless law, each $g_i$ is therefore an exact member of
\noderef{BoundaryEdges}, colored either $a$ or $b$.  Closedness in
\noderef{TreeColorSetClosed}, through the exact filter
\noderef{BoundaryColorEdges}, excludes color $a$.  Thus $g_1,g_3$ are
distinct $b$-colored edges incident with the same vertex.  They are
adjacent by \noderef{AdjacentEdges}, contradicting properness.  Hence
$P_1$ and $P_3$ are support-disjoint after all.

Support disjointness gives $a_3\ne a_1$ and $b_3\ne b_1$, while full
component disjointness gives $a_3\ne a_2$ and $b_3\ne b_2$.  Finally
apply \noderef{ClosedColorExitPathBoundaryEdge} to $P_3$.  It supplies an
exact edge $f_3$ in the support of $P_3$, incident with $b_3$, and lying
in the $b$-colored boundary of $X$.  These are all the required
conclusions.
\end{proof}
